\section{The ordinary-$abc$ quartics}
The 68-variable source uses two products to test four field quantities. In its notation let $C=s-S$, $I=t-T_H$, $J=\beta^2 f(H)f(U)$ and $K'=\beta^2 f(H-R_0)f(V)$. With $q_j=1/w_j$, $c=q_0+q_1$, the single replacement test is
\[
 W=1+c C^4+c^2J^2-c^3\bar I^4-c^4\bar K'{}^2 .
\]
Set $D\Tr(W)=\Nm(W-1)$ and $2\tau_j=1+q_{1-j}^2+q_jD^2$. The large pole predicate receives $b_j=w_j\tau_j^2$. At a consumed prime the pole inequality gives $v(D)>\mu/4$ and $v(q_{1-j})>-\mu/4$, hence $v(c)=-\mu$. If $W$ had poles in both slots the norm-over-trace estimate would give $v(D)<0$. Thus $v(W)\ge0$ in some slot. The four nonconstant terms have pairwise distinct valuation classes modulo four since $J,K'$ have even valuation. The last term can tie with the constant at valuation zero. Any negative minimum would be unique, so all four terms are integral in this slot. The strictly negative valuation of $c$ implies all four raw tests have positive valuation. This gives the source input congruence.

For completeness one avoids the finitely many primes dividing $1+q_{1-j}^2$ at the translated contacts once the shifts are fixed. At the prepared successful slot $W\equiv1$. At the opposite slot the two reciprocal products are units by the v65 prime avoidance, so $W$ has a pole by uniqueness of a negative minimum. Thus $v(D)>0$. At $X=0$ the source chart makes all four tests zero, so $W=1$ and $D=0$. The same disjoint prime choice supplies both prepared pole inputs. The proof and independent symbolic checker in the bundle give the real, two-adic and denominator conditions for the large pole predicate.

Use $(C^2\pm c\bar I^2)$ and $(J\pm c\bar K')$ as pairs of field coordinates. The resulting 62-coordinate system has 50 quadrics. Their squared sum has 830 monomials and exact height $B^2+4091$ for every odd $B\ge81$. On the cofinal shifts $B=n^4$ a diagonal rescaling keeps all 830 monomials and gives exact height $B+4091$ for every odd $n\ge3$. Four further invertible row replacements and positive row weights give 757 monomials at exact height $B+758541$ on the same subfamily.

At 68 variables and odd fourth-power shifts, an invertible change of the reference coordinates, row operations, positive integer row weights, and diagonal scalings give a primitive quartic with 693 monomials. Its exact height is $\max(B+758542,2B)$ for odd $n\ge7$, hence $2B$ for all odd $n\ge31$. The source has 1,078 monomials and height $B^2+1$. These identities concern coefficients. A sufficient defining shift under ordinary rational $abc$ is still ineffective.

\section{Degree bounds at fixed variable counts}
The following are upper bounds for unexpanded integer determinants. The first two rows use the source arithmetic premises. The next two add the indicated versions of weak $abc$. The final row adds ordinary rational $abc$.
\begin{center}\small
\begin{tabular}{llrr}
\toprule
&Variables&Source degree bound&New degree bound\\
\midrule
General&25&4,863,007,848,923,136&238,739,783,680\\
Fully specified, $L=20483$&27&99,279,020,543,180,800&1,941,079,851,008\\
Weak $abc$ at $121/100$&13&43,211,777&2,525,697\\
Weak $abc$ at $16/15$&13&39,858,177&2,206,209\\
Ordinary rational $abc$&11&230,932,097&36,482,433\\
\bottomrule
\end{tabular}\end{center}

In these finite-algebra sources, each detector product has factors $1+a_i$ whose nonzero $a_i$ have odd valuation at the good places used for local extraction. Replace the product by $W_+=\prod_i(1+a_i^2)$. At those places $v(1+a_i^2)=2\min(0,v(a_i))$. This gives the original local inequalities, and $W_+$ has residue one on the successful prepared slot. On any rational point character with the raw denominators defined, $W_+$ is strictly positive at both real embeddings. Its rational component and field norm are therefore nonzero. The source imposes corresponding nonvanishing conditions by norms across a large finite algebra. We can omit precisely those numerator atoms forced by positivity after a rational character is selected. The denominators of the raw tests, product and field norm remain protected. This order avoids presuming positivity on a complex character. The bundle gives the finite-place, real and denominator checks for each row.

For the 11-variable row, positive denominators can be used one step earlier as well. On a test $x$ whose $ux$ has odd valuation, replace $[x/(1+ux)]^2$ by $x^2/(1+(ux)^2)$. They have the same valuation at each consumed place. The latter denominator is positive in both real embeddings on a rational point character. This removes further factors from the scalar guard, after protecting the denominators defining $x$.

A second reduction concerns the final pole generators $u_0,\ldots,u_{s-1}$, $s\le3$. Over any fixed character of the preceding point algebra, their remaining choices are signs of square roots, and zero roots are allowed by the pole obstruction. For $s\ge2$ the single expression $E=\sum_{i<s}2^i u_i^3$ distinguishes all distinct sign tuples. Indeed embed all conjugates at a rational base specialization in a finite Galois extension of $2$-power degree, obtained by adjoining all conjugate square roots at each layer. At a prime over two, use the integer-normalized valuation with $e=v(2)$ a power of two. The summands of a nonzero sign difference, after removing the common factor two, have valuations $ie+3v(u_i)$ in distinct classes modulo three. The least valuation is unique. For $s=1$ use $E=u_0$.

To separate earlier point choices one uses the nonzero point generators $t_i$ and the source scalar polynomial bounds $M_i\ge |t_i|$ on all complex characters. With another positive envelope $B'\ge\max(1,|E|)$ set $\phi_{m+1}=E/(4B')$ and recursively $\phi_i=t_i(1+\phi_{i+1})/(4M_i)$. At the first different point sign the two resulting discs are disjoint, since $|\phi_{i+1}|\le1/3$. Every common preceding map is injective. Clearing these positive scalar denominators gives an integral separating element $\Omega$ with nonzero character values. If $G$ is the product of the retained scalar norm guards, use $\det(I-YG M_\Omega)$, with the original input factor if required. A rational root gives a rational character since its eigenvalue is rational and distinguishes that character from every conjugate. Eigenvalues of multiplication are the character values even for a nonreduced finite algebra. Conversely a rational character with valid guards supplies $Y=1/(G\Omega)$. The explicit 27-variable member has coefficient height below $2^{2^{54}}$. Source graphs, separator formulas and degree ledgers are included.

For two poles one can improve the separator using a rational function. If $x^2=d$, $z^2=e$ are rational and $|x|,|z|\le1/2$, then $x^3+z/(1+z)$ is rational only if both $x,z$ are rational. Its possibly irrational part is $dx+z/(1-e)$. If exactly one root is irrational this cannot be rational. If both are irrational the relation forces $z=tx$ with $t\in\Q$ and $d^2t^2-t-d=0$. Hence $1+4d^3$ is a square, whereas the only affine rational points of $y^2+y=d^3$ have $d=0$. This last result also follows from Fermat's theorem at exponent three, as detailed in the bundle. Normalize each pole root by twice its positive envelope to obtain $x,z$, and use $(x^3+z/(1+z))/8$ for the final digit. A rational spectral value first forces all preceding point coordinates to be rational by the disjoint-disc argument. The lemma then applies in their rational fiber. Writing the resulting fraction as $P/Q$, the integral matrix pencil is $M_Q-YG M_P$. All possible denominator values are nonzero. This construction gives the displayed 25-variable and two 13-variable rows.
